\phantomsection
\section*{KẾT LUẬN}
\addcontentsline{toc}{section}{KẾT LUẬN}
\label{sec:ketluan}

Bài tập lớn đã xây dựng cơ sở dữ liệu quản lý cư dân và biến động cư trú tại một phường trên PostgreSQL 17. Mô hình \textbf{lưu riêng các đợt cư trú theo thời gian}: mỗi đợt của một người tại một hộ có ngày bắt đầu bắt buộc và ngày kết thúc có thể để trống; mỗi biến động tham chiếu đợt cũ, đợt mới hoặc cả hai tùy loại sự kiện, đồng thời ghi tài khoản và thời điểm nhập. Thiết kế hỗ trợ tra cứu lịch sử đã ghi nhận, còn việc bảo vệ các bản ghi phụ thuộc quyền và trigger. Kết quả cụ thể của nhóm như sau.

\begin{enumerate}[leftmargin=*]
    \item \textbf{Thiết kế dữ liệu.} Từ bài toán và 14 chức năng cấp 2, nhóm xây dựng mô hình thực thể liên kết theo ký hiệu Chen rồi ánh xạ sang 7 bảng với 35 cột, 7 khóa chính, 7 khóa ngoại và 5 ràng buộc UNIQUE; hai UNIQUE trên cột đợt cũ và đợt mới thể hiện hai liên kết 1 -- 1 giữa đợt cư trú và biến động. Thực thể Quá trình cư trú giải quan hệ nhiều -- nhiều giữa cư dân và hộ, đồng thời là nguồn dữ liệu tham chiếu để tra cứu các đợt đã lưu theo mã cư dân. Kết quả phân tích BCNF theo tập phụ thuộc hàm đã xét được nêu ở mục \ref{subsec:chuanhoa}. Ngày biến động trùng với ngày của đợt là dữ liệu lặp giữa hai bảng, được trigger B09 giữ cho khớp.
    \item \textbf{Ràng buộc nghiệp vụ.} Ngoài khóa và CHECK, các quy tắc B06 -- B15 được cài bằng ràng buộc loại trừ GiST (không chồng lấn đợt cùng loại), trigger kiểm tra hình dạng và ngày của biến động, trigger giữ địa chỉ đã có lịch sử, ba hàm chuyển nội bộ, chuyển đi và chuyển đến thực hiện trong một giao dịch, và trigger kiểm tra ngày hợp lý. Cơ sở dữ liệu dựng mới nạp được bộ dữ liệu kiểm thử gồm 636 bản ghi, bước kiểm tra dữ liệu không phát hiện vi phạm, và các câu lệnh cố tình vi phạm đều bị từ chối với đúng mã lỗi.
    \item \textbf{Truy vấn.} Báo cáo trình bày 10 câu đại số quan hệ, gồm đủ phép chọn, chiếu, kết nối theo điều kiện, đổi tên, hợp, giao, hiệu, chia và phép gom nhóm của đại số quan hệ mở rộng, cùng 13 câu SQL dùng kết nối nhiều bảng, kết nối ngoài, gom nhóm có HAVING, truy vấn con tương quan, truy vấn con trong FROM, CTE, hàm cửa sổ RANK() và phép chia viết bằng hai lớp NOT EXISTS. Kết quả của cả 13 câu QSQL khớp với dữ liệu đối chiếu khi chạy lại trên cơ sở dữ liệu dựng mới. Các biểu thức đại số quan hệ được kiểm chứng bằng những truy vấn SQL đối chứng tương ứng.
    \item \textbf{Bảo mật.} Hai ví dụ kiểm thử quyền dùng phân quyền theo nguyên tắc quyền tối thiểu (role chỉ đọc bị từ chối khi sửa dữ liệu, role quản trị không xóa được lịch sử) và khung nhìn chỉ công bố số liệu tổng hợp theo tổ. Cả hai đã được kiểm tra bằng phiên đăng nhập thật của từng role và bằng cách kiểm tra quyền đã cấp. Ví dụ SQL Injection minh họa vì sao phải truyền dữ liệu người dùng như tham số thay vì ghép chuỗi vào câu lệnh động.
\end{enumerate}

Cột vai trò của tài khoản chỉ là nhãn nghiệp vụ, không ngăn tài khoản thống kê ghi biến động (X01). B15 chưa kiểm tra được ngày sinh hay ngày lập hộ ở tương lai khi chưa có đợt cư trú (X06). B10 chỉ được bảo đảm khi ghi biến động qua ba hàm chuyển; role quản trị vẫn sửa trực tiếp được bảng đợt cư trú, nên vẫn có thể đóng một đợt mà không ghi biến động. Role quản trị vẫn có quyền sửa hai bảng lịch sử trong phạm vi trigger cho phép và có quyền đọc, thêm, sửa và xóa bản ghi trên bảng tài khoản, nên có thể đổi người ghi của biến động cũ, đổi tên hiển thị hoặc mở khóa tài khoản. Đồ án cũng chưa triển khai bảo mật mức dòng, nhật ký kiểm toán đầy đủ và mã hóa kết nối.

Về nghiệp vụ, mô hình còn đơn giản hơn thực tế quản lý cư trú ở bốn điểm. Thứ nhất, cư dân chưa có số định danh cá nhân (số CCCD) và họ tên được phép trùng, nên không phân biệt được hai người trùng cả họ tên lẫn ngày sinh ngoài mã cư dân do hệ thống cấp. Thứ hai, hộ chưa có chủ hộ và chưa lưu quan hệ của từng thành viên với chủ hộ. Thứ ba, việc đổi từ tạm trú sang thường trú trong cùng một hộ chưa được ghi thành biến động, vì CHUYEN\_NOI\_BO bắt buộc hai hộ khác nhau (B08). Thứ tư, chưa có loại biến động cho khai sinh và khai tử; hiện tại người mới sinh chỉ có thể được ghi như chuyển đến, còn người mất như chuyển đi. Nhóm triển khai cơ sở dữ liệu trong môi trường kiểm thử và dùng bộ dữ liệu tự tạo để ghi nhận các kết quả trong báo cáo.

Trong thời gian tới, cơ sở dữ liệu có thể được hoàn thiện theo các hướng: bổ sung số định danh cá nhân có ràng buộc UNIQUE cho cư dân, chủ hộ và quan hệ với chủ hộ, cùng các loại biến động đổi loại cư trú, khai sinh và khai tử; chỉ cho ghi biến động qua các hàm chuyển; thu hẹp quyền của role quản trị trên bảng tài khoản và cột người ghi của biến động, bổ sung bảng nhật ký kiểm toán ghi lại mọi lệnh sửa dữ liệu lịch sử; ràng buộc vai trò được phép ghi biến động nếu nghiệp vụ yêu cầu; áp dụng bảo mật mức dòng để mỗi cán bộ chỉ xem dữ liệu của tổ mình; và bổ sung chức năng thống kê lịch sử theo nhu cầu tra cứu.
